\section{Methodology}
Two kinds of numerical investigations were undertaken: a stationary airfoil at different angles of attack, and a pitching airfoil following a prescribed rotational motion $\alpha(t)$. The finite-volume code OpenFOAM, specifically the Foundation distribution, was used to set up both kinds of simulation \cite{weller}. 
%\todo{confirm the version for each campaign: the original methodology specifies OpenFOAM v13, while some case records identify a development build.}

\subsection{Static NACA0012 Airfoil}
\subsubsection{Mesh}
For both the stationary and pitching NACA0012 airfoils, the computational domain is a disk of radius $30c$. The mesh is discretized using transfinite interpolation, with quadrilateral elements throughout the two-dimensional plane. An overview and close-up are provided in Figure~\ref{fig:mesh}.

\fig{figs/mesh_figure.pdf}{0.95}{Computational mesh: the circular domain and near-field refinement, with close-ups of the leading edge and trailing edge.}{fig:mesh}

To keep the boundary layer sufficiently resolved, the first-cell height is taken as a function of the Reynolds number,
\begin{equation}
\frac{h_1}{c}=\frac{0.05}{\sqrt{Re}},
\end{equation}
which is approximately one hundredth of the Blasius boundary-layer thickness for a flat plate at $x=c$. The mesh quantities are summarized in Table~\ref{tab:mesh}. The flat-plate estimate provides a mesh-design scale rather than a measurement of the separated airfoil boundary layer.

\begin{table}[htbp]\centering\small
\caption{Mesh quantities for the $Re=500$ NACA0012 airfoil.}\label{tab:mesh}
\begin{tabularx}{\linewidth}{@{}>{\raggedright\arraybackslash}p{.57\linewidth}X@{}}
\toprule Quantity & Value \\
\midrule
Total cells & 63\,360 \\
Cells on each surface / across the trailing-edge base / radial & 160 / 32 / 180 \\
First-cell height & $2.24\times10^{-3}c$ \\
Radial growth ratio & 1.035 \\
Cells inside the estimated boundary layer ($\delta_{99}=0.224c$ at $x=c$) & 44--49 per wall-normal line \\
Maximum non-orthogonality / skewness / aspect ratio & $64^\circ$ / 0.82 / 7.5 \\
\bottomrule
\end{tabularx}
\end{table}

For the stationary cases, the angle of attack $\alpha$ is changed by rotating the freestream direction. For pitching cases, the whole mesh is rotated rigidly around a pivot.

\subsubsection{Solver Setup}
As the flow regime is laminar, the shear layer near the no-slip wall is an important small scale to resolve. The governing equations were solved directly; no turbulence models were employed. The solver settings are given in Table~\ref{tab:static_solver}.

\begin{table}[htbp]\centering\small
\caption{Solver setup for the time-accurate stationary-airfoil simulations.}\label{tab:static_solver}
\begin{tabularx}{\linewidth}{@{}>{\raggedright\arraybackslash}p{.28\linewidth}X@{}}
\toprule Item & Setting \\
\midrule
Pressure--velocity coupling & PISO (PIMPLE, one outer corrector), two pressure correctors, one non-orthogonal corrector \\
Time derivative & Second-order implicit backward \\
Convection & Second-order linear upwind \\
Gradient, Laplacian & Second-order central (Gauss linear, non-orthogonal correction) \\
Linear solvers & $p$: GAMG, tolerance $10^{-8}$ on the final corrector; $\mathbf{U}$: symmetric Gauss--Seidel, tolerance $10^{-9}$ \\
Boundary conditions & Airfoil: no-slip; far field: freestream velocity and pressure; span: two-dimensional (\texttt{empty}) \\
\bottomrule
\end{tabularx}
\end{table}

\subsubsection{Time Stepping}
The time step is fixed within each run stage; adaptive time stepping is disabled. The prescribed run procedure has three phases:
\begin{enumerate}
\item \textbf{Start-up.} An impulsive start at $\Delta t=4\times10^{-5}$ up to $t=0.5$.
\item \textbf{Transient.} The step is chosen from the Courant number measured at the end of start-up, with a nominal cap near $\Delta t=6\times10^{-4}$. The run advances until the lift history reaches a limit cycle: over two consecutive windows, its mean, amplitude, and frequency must agree to 1\%.
\item \textbf{Snapshot record.} The target is twenty limit-cycle periods, with 32 snapshots per shedding period.
\end{enumerate}
Reported production time steps were $\Delta t=4.3$--$6.3\times10^{-4}$, and the maximum Courant number stayed below 0.8. The upper reported step slightly exceeds the nominal cap in the original procedure. Actual averaging windows are recorded separately; the target recording duration is not assumed to have been achieved in every transferred case. Time and time step are nondimensionalized by $c/U_\infty$.

\FloatBarrier
\subsection{Pitching NACA0012}
\subsubsection{Airfoil Motion Using the Arbitrary Lagrangian--Eulerian Formulation}
The same mesh topology as the static case is reused for the pitching airfoil. The circular symmetry of the O-grid is exploited to implement pitching motion. The entire mesh rotates relative to the flow, and time derivatives in the moving mesh are related to those in the stationary coordinate frame through the arbitrary Lagrangian--Eulerian (ALE) transformation:
\begin{align}
\left.\frac{\partial\phi}{\partial t}\right|_{\boldsymbol\xi}
&=\left.\frac{\partial\phi}{\partial t}\right|_{\mathbf{x}}
+\sum_i\frac{\partial\phi}{\partial x_i}
\left.\frac{\partial x_i}{\partial t}\right|_{\boldsymbol\xi}\\
&=\left.\frac{\partial\phi}{\partial t}\right|_{\mathbf{x}}
+\mathbf{u}_m\cdot\nabla\phi.
\end{align}
Here $\mathbf{u}_m$ is the mesh velocity. For rigid-body rotation, it is given by
\begin{equation}
\mathbf{u}_m=\boldsymbol\Omega\times\mathbf{r},
\end{equation}
where $\mathbf{r}$ is measured from the pivot. For the two-dimensional case,
\begin{equation}
\boldsymbol\Omega(t)=\dot\alpha(t)\,\hat{\mathbf{k}},
\end{equation}
where $\alpha(t)$ is the prescribed rotational displacement, with its sign defined consistently with mesh rotation. This transformation changes the material derivative $D\mathbf{u}/Dt$ and hence the Navier--Stokes momentum equation as follows:
\begin{equation}
\left.\frac{\partial\mathbf{u}}{\partial t}\right|_{\boldsymbol\xi}
+[(\mathbf{u}-\mathbf{u}_m)\cdot\nabla]\mathbf{u}
=-\nabla p_k+\nu\nabla^2\mathbf{u},
\qquad \nabla\cdot\mathbf{u}=0,
\end{equation}
where $p_k=p/\rho$ is kinematic pressure and $\nu$ is kinematic viscosity.

\subsubsection{Solver Setup}
In the prescribed-step setup, each unsteady simulation is started from a converged steady-flow solution at $\alpha=0$, obtained using the steady-state SIMPLE algorithm. The transient PIMPLE algorithm is then used to calculate the time-evolving pressure and velocity fields. The settings recorded in the original methodology are retained in Table~\ref{tab:pitch_solver}.

\begin{table}[htbp]\centering\small
\caption{Pitching setup recorded for the prescribed-step NACA0012 campaign.}\label{tab:pitch_solver}
\begin{tabularx}{\linewidth}{@{}>{\raggedright\arraybackslash}p{.28\linewidth}X@{}}
\toprule Item & Setting \\
\midrule
Pressure--velocity coupling & PIMPLE: two outer correctors, two pressure correctors, two non-orthogonal correctors; flux corrected for mesh motion \\
Time step & Fixed, $\Delta t=5\times10^{-4}$ \\
Time derivative / convection / diffusion & Second-order backward / linear upwind / central (Gauss linear, non-orthogonal correction) \\
Linear solvers & $p$: GAMG, tolerance $10^{-7}$ on the final corrector; $\mathbf{U}$: symmetric Gauss--Seidel, tolerance $10^{-6}$ \\
Boundary conditions & Airfoil: no-slip relative to the moving wall; far field: freestream velocity and pressure; span: two-dimensional (\texttt{empty}) \\
Output & Forces every time step; flow fields every 0.05 (every 0.0125 during the $45^\circ$ ramp) \\
\bottomrule
\end{tabularx}
\end{table}

%\todo{reconcile the final-corrector tolerances in this original table with the final archived case dictionaries.} These settings describe the earlier prescribed-step campaign, not every pitching result. The current $Re=500$, $30^\circ$ case uses a quarter-chord pivot; its dictionary specifies one outer corrector, two pressure correctors, one non-orthogonal corrector, and $\Delta t=2\times10^{-4}$. The separate $Re=1000$ flat-plate benchmark uses a leading-edge pivot \cite{eldredge2010}.

\FloatBarrier
\subsubsection{Input Function}
For the pitch-up input function in the prescribed-step campaign, a continuously differentiable approximation to a step function is used:
\begin{equation}
\alpha(t)=\Delta\alpha\,s^3(6s^2-15s+10),
\qquad s=\min\!\left(\max\!\left(\frac{t}{T},0\right),1\right).
\end{equation}
This $\alpha(t)$ has $\dot\alpha=\ddot\alpha=0$ at the beginning and end of the maneuver, which avoids the singularity associated with an instantaneous step and improves numerical regularity. Here $T$ is the ramp duration and $\Delta\alpha$ is the prescribed change in angle.

The later $Re=500$ maneuver uses a separate Eldredge-type smoothed pitch-up-and-hold input, with nominal reduced rate $K=\dot\alpha_0c/(2U_\infty)=0.4$ and final angle $30^\circ$. Its motion table is clamped at the maximum angle for the hold. The corresponding angle, angular velocity, and angular acceleration are shown in Figure~\ref{fig:kinematics}; this input is distinguished from the polynomial step approximation above.

\fig{assets/pitch_kinematics.png}{0.58}{Angle, angular velocity, and angular acceleration for the later $Re=500$ NACA0012 pitch-up-and-hold input.}{fig:kinematics}
\FloatBarrier
